\documentclass[11pt,letterpaper]{article}
\PassOptionsToPackage{hyphens}{url}
\usepackage[margin=.83in,top=.8in,bottom=.8in,headheight=15pt]{geometry}
\usepackage[T1]{fontenc}
\usepackage{lmodern,amsmath,amssymb,mathtools,bm,booktabs,array,graphicx,tikz}
\usetikzlibrary{arrows.meta,positioning,calc,decorations.pathreplacing,patterns}
\definecolor{ink}{HTML}{17344D}\definecolor{accent}{HTML}{087E8B}\definecolor{pale}{HTML}{EDF5F5}\definecolor{gold}{HTML}{AD6B19}
\providecommand{\dd}{\,\mathrm{d}}\providecommand{\R}{\mathbb{R}}
\providecommand{\E}{\mathbb{E}}\providecommand{\Var}{\operatorname{Var}}
\providecommand{\Cov}{\operatorname{Cov}}\providecommand{\diag}{\operatorname{diag}}
\providecommand{\norm}[1]{\left\lVert#1\right\rVert}
\providecommand{\abs}[1]{\left|#1\right|}
\usepackage{microtype,hyperref,fancyhdr,enumitem,amsthm,sectsty,needspace}
\allsectionsfont{\raggedright}
\hypersetup{colorlinks=true,linkcolor=ink,urlcolor=accent,pdfauthor={Zan Ahmad},pdftitle={Fetal pathways and signed junction balances / Review}}
\setlength{\parindent}{0pt}\setlength{\parskip}{6pt}\linespread{1.07}
\setlist{itemsep=3pt,topsep=4pt,leftmargin=*}
\emergencystretch=2em\clubpenalty=10000\widowpenalty=10000
\pagestyle{fancy}\fancyhf{}\fancyhead[L]{\small Mathematical Cardiology}\fancyhead[R]{\small Lecture 7}
\fancyfoot[L]{\footnotesize Zan Ahmad}\fancyfoot[R]{\footnotesize\thepage}\renewcommand{\headrulewidth}{.3pt}
\newtheorem{proposition}{Proposition}\newtheorem{theorem}{Theorem}\newtheorem{lemma}{Lemma}
\theoremstyle{definition}\newtheorem{definition}{Definition}\newtheorem{example}{Example}
\numberwithin{equation}{section}
\begin{document}
\section*{Fetal pathways and signed junction balances --- review sheet}
Use this after attempting the worked examples and practice. Explain each claim in words before calculating.

\section*{Key ideas}
\begin{itemize}
\item The placenta supplies fetal exchange; pulmonary perfusion is not zero.
\item Foramen ovale links atria; ductus arteriosus links pulmonary artery and aorta.
\item Flow reversal and pathway closure are different.
\item At steady junctions, inflow equals outflow with the chosen signs.
\item A topology change requires a newly consistent state.
\end{itemize}
\section*{Common errors and repairs}
\textbf{Fetal blood is blue before it reaches the lungs.} Blood is red; oxygen content varies, and the placenta is the main fetal exchange organ.

\textbf{Zero flow proves permanent closure.} Equal pressures or directional restriction can also produce zero flow.

\textbf{All other flows stay fixed when a shunt is erased.} Recheck every junction balance after changing the connection rule.

\section*{Retrieve without your notes}
\begin{itemize}
\item Which connection bypasses part of hepatic circulation?
\item What does Qd < 0 mean?
\item What do the current equations fail to predict?
\end{itemize}
\section*{Readiness check}
\begin{itemize}
\item Trace placental supply, foramen ovale and ductus arteriosus pathways using named anatomy.
\item Distinguish pressure-driven flow reversal from functional or anatomical closure.
\item Calculate a signed ductus flow and a one-way foramen surrogate.
\item Check four junction balances before and after a prescribed topology change.
\item State which birth-transition mechanisms the network calculation leaves out.
\end{itemize}
\end{document}
